% !PRINTED 303
\weberpage{303}

Wir können aber auch noch die Voraussetzung aufgeben, dass für $\alpha$ und
$\beta$ keine der Functionen (1) verschwindet, wenn $\alpha$ und $\beta$
nicht unter den Wurzeln von $f(x)$ vorkommen.

Denn dann ist unser Theorem anwendbar auf zwei Werthe, von denen der erste
etwas grösser als $\alpha$, der andere etwas kleiner als $\beta$ ist, und
die dieselben Wurzeln einschliessen, wie $\alpha$ und $\beta$; d.\,h. es
ist, wenn wir mit $\mathfrak{z}$ die Anzahl dieser Wurzeln bezeichnen
\[
V(\alpha + 0) - V(\beta - 0) = \mathfrak{z} + 2h,
\]
und da
\[
V(\alpha + 0) = V(\alpha), \quad V(\beta - 0) = V(\beta) + 2h',
\]
so folgt
\[
V(\alpha) - V(\beta) = \mathfrak{z} + 2(h + h'),
\]
was wieder der Ausdruck unseres Theorems ist.

Wenn aber $\alpha$ oder $\beta$ selbst zu den Wurzeln von $f(x)$ gehört,
dann ist unser Theorem nur dann richtig, \gesperrt{wenn diese Grenzwerthe nicht mit
zum Intervall gezählt werden}.

Das \gesperrt{Budan-Fourier'sche} Theorem giebt zwar nicht, wie der
\gesperrt{Sturm'sche} Satz, eine vollständig sichere Entscheidung über die
Zahl der Wurzeln in einem Intervall, es kann aber doch in manchen Fällen
den \gesperrt{Sturm'schen} Satz ersetzen; denn hat man das Intervall
$(\alpha, \beta)$ so weit eingeschränkt, dass kein oder nur ein
Zeichenwechsel von $\alpha$ bis $\beta$ verloren geht, so folgt mit
Sicherheit, dass im ersten Falle keine, im zweiten eine und nur eine
Wurzel im Intervall liegt.

Die Abgrenzung der Wurzeln kann also durch die Budan'sche Reihe nur dann
vollständig gegeben werden, wenn nicht beim Durchgang durch einen Werth
$\xi$ gleichzeitig mehrere Zeichenwechsel verloren gehen. Dies Verhalten
kann, wie klein auch das Intervall $(\alpha, \beta)$ gewählt sein mag,
nicht mit Sicherheit erkannt werden, und man wird also, wenn nach einer
angemessenen Einengung des Intervalles die Abgrenzung nicht gelungen ist,
doch zu einem anderen Verfahren, in letzter Instanz zum
\gesperrt{Sturm'schen} Satze greifen müssen.
